% ====================================================================
% Section 2 (v2): Related work.
% Positions the contribution as parsimony, not SOTA competition.
% ====================================================================
\section{Related Work}
\label{sec:related}

\paragraph{Engineered load balancing.}
Modern multipath load balancers spread traffic to avoid hot spots.
ECMP~\cite{ecmp} splits flows across equal-cost shortest paths by hashing, with
no congestion awareness. CONGA~\cite{conga} adds that awareness: it enumerates a
set of candidate paths and selects among them using in-network congestion
feedback, rerouting at flowlet granularity. DRILL~\cite{drill} pushes the
decision to the switch and the microsecond, making per-packet local choices
based on queue occupancy with minimal state. These systems define the operating
points we measure against. Our aim is not to outperform them---on the
topologies where they are strong, they remain strong---but to ask how much of
their behaviour is recoverable from a far simpler, physically-motivated rule. Transport-layer multipath---MPTCP~\cite{ford2013mptcp} and the in-progress multipath extension of QUIC~\cite{ietfquicmultipath}---is orthogonal to this in-network family: endpoints spread load across whole paths end to end, while the balancers above, and our core, choose next hops inside the network.

\paragraph{Physics- and nature-inspired routing.}
Using physical or biological analogies for routing has a long history:
ant-colony optimisation~\cite{dicaro1998antnet}, electrical-network and resistance analogies,
potential- and field-based forwarding, and physarum-inspired path formation~\cite{tero2010physarum}.
These approaches typically introduce a mechanism and demonstrate that it works.
Our contribution differs in emphasis in two ways. First, we run the analogy in
reverse as an ablation: rather than showing that added physics helps, we add a
great deal of physics and then measure how much can be removed, arriving at a
two-term residue. Second, we frame the outcome as a statistical equivalence
against strong engineered baselines, rather than as a demonstration that the
physical method is novel or superior.

\paragraph{Scope of comparison.} We benchmark against CONGA and DRILL because
they are the state-of-the-art in adaptive, congestion-aware load balancing---the
exact problem the physical core addresses. Schemes that do not solve that
problem are out of scope, not because they are unimportant but because they
answer a different question: shortest-path and ECMP appear only as blind
references (they ignore congestion by construction), and legacy access or
medium-arbitration protocols (token passing, spanning-tree-era schemes) are
not load balancers at all. Equivalence with the best congestion-aware
balancers is the claim that carries weight; equivalence with a non-adaptive
scheme would not.


\paragraph{Backpressure and queue-gradient methods.}
Backpressure scheduling~\cite{tassiulas1992stability} and its delay-aware
refinements~\cite{ji2013delay} route along gradients of per-queue backlog,
and are throughput-optimal under stationary arrivals. The connection to the
present work is direct: our load term is itself a queue-gradient signal, and
the two-term core can be read as a backpressure-like rule augmented with a
shortest-path bias and a slow scar memory. The distinction we draw is not
mechanistic but evaluative---we ask how close such a rule sits to engineered
multipath balancers, and report the answer as an equivalence with a stated
margin rather than as a throughput-optimality proof.

\paragraph{Learning-based routers.}
A parallel line replaces hand-built rules with learned policies, from
Q-routing~\cite{boyan1994packet} to graph-aware deep reinforcement
learning~\cite{valadarsky2017learning}; recent surveys document a growing
body of such routers~\cite{xiao2021survey}. These trade interpretability for
adaptivity. Our contribution is complementary and deliberately opposite in
spirit: a fixed, inspectable rule with no training, offered not as a
performance ceiling but as a parsimony baseline against which the value added
by learning---or by engineered path enumeration---can be measured.

\paragraph{Per-packet and per-flow congestion state.}
A related family carries congestion state with the traffic itself. ECN
~\cite{ramakrishnan2001ecn} marks packets that traversed a congested queue;
DCTCP~\cite{alizadeh2010dctcp} reacts to the fraction marked; ConEx
~\cite{briscoe2014conex} re-exposes that signal for accountability.
These mechanisms share with our scar term the idea that a link should carry a
memory of recent trouble. The difference is where the memory lives and acts:
these feed endpoint congestion control, leaving the in-network forwarding
decision stateless, whereas our scar modulates the next routing choice locally,
without a round trip. The flaky-link benchmark (Section~\ref{sec:freshness})
is where that distinction earns its keep.

\paragraph{Equivalence testing.}
Most networking evaluations test for a difference (one method beats another).
We instead test for \emph{equivalence}, using TOST (two one-sided
tests), which asks whether two methods are indistinguishable within a stated
margin. This is the statistically appropriate tool when the claim is ``as good
as'', not ``better than'', and it guards against the familiar error of reading a
non-significant difference as evidence of equivalence. To our knowledge,
equivalence testing is uncommon in the routing-evaluation literature, and its
use here is deliberate: the entire thesis is one of equivalence with parsimony.

\paragraph{Applicability conditions.}
A recurring weakness of method papers is silence about when the method fails. We
attach to our result an explicit, topology-only predictor (the tube/sp ratio,
Section~\ref{sec:applicability}) that says in advance where the equivalence
holds and where engineered routing retains an edge. This places the
contribution closer to the spirit of complex-systems and network-science work,
where structural predictors of dynamical behaviour are themselves the object of
study, than to a conventional systems-performance paper.


\medskip
Conceptually adjacent, though router-free, is the literature on dynamic
averaging and load-balancing processes on graphs, e.g.\ Alistarh et
al.~\cite{alistarh2022dynamic}, where iterative local exchanges converge
to equilibria fixed by graph structure. The attractor reported in this
paper is the routing-level analogue of that structural determinism,
observed empirically across heterogeneous mechanisms rather than proved
for a single dynamics.
